\documentclass[10pt,a4paper]{article}
\usepackage[utf8]{inputenc}
\usepackage[T1]{fontenc}
\usepackage[provide=*,spanish]{babel}
\usepackage{lmodern,microtype}
\usepackage[a4paper,left=18mm,right=18mm,top=17mm,bottom=20mm]{geometry}
\usepackage{amsmath,amssymb,booktabs,array,longtable,tabularx}
\usepackage[table]{xcolor}
\usepackage{fancyhdr,titlesec,enumitem,listings}
\usepackage[most]{tcolorbox}
\usepackage[hidelinks]{hyperref}

\definecolor{navy}{HTML}{17324D}
\definecolor{blue}{HTML}{2D6A9F}
\definecolor{pale}{HTML}{EAF2F8}
\definecolor{light}{HTML}{F5F7FA}
\definecolor{graytext}{HTML}{5D6D7E}
\definecolor{bodytext}{HTML}{202830}
\color{bodytext}
\setlength{\parindent}{0pt}
\setlength{\parskip}{5pt}
\setlength{\emergencystretch}{1em}
\setlist{nosep,leftmargin=5mm}
\renewcommand{\arraystretch}{1.22}
\titleformat{\section}{\Large\bfseries\color{navy}}{\thesection.}{0.45em}{}
\titleformat{\subsection}{\large\bfseries\color{blue}}{\thesubsection.}{0.45em}{}
\titlespacing*{\section}{0pt}{12pt}{7pt}
\titlespacing*{\subsection}{0pt}{9pt}{4pt}
\pagestyle{fancy}
\fancyhf{}
\renewcommand{\headrulewidth}{0pt}
\renewcommand{\footrulewidth}{0.35pt}
\fancyfoot[L]{\footnotesize\color{graytext}Firmware PixiePico -- Nota técnica PLL/PIO}
\fancyfoot[R]{\footnotesize\color{graytext}Página \thepage}
\newcolumntype{Y}{>{\raggedright\arraybackslash}X}
\newcolumntype{R}[1]{>{\raggedleft\arraybackslash}p{#1}}
\newtcolorbox{technicalbox}{
  colback=pale,colframe=blue,boxrule=0.55pt,arc=0pt,
  left=2mm,right=2mm,top=1.4mm,bottom=1.4mm}
\lstdefinestyle{code}{
  basicstyle=\ttfamily\scriptsize,backgroundcolor=\color{light},
  frame=single,rulecolor=\color{graytext!35},breaklines=true,
  columns=fullflexible,keepspaces=true,showstringspaces=false,
  xleftmargin=1mm,xrightmargin=1mm}
% Fallback autocontenido para no depender de siunitx.
\newcommand{\num}[1]{#1}
\newcommand{\SI}[2]{#1\,\ensuremath{\mathrm{#2}}}
\newcommand{\mega}{M}
\newcommand{\hertz}{Hz}

\begin{document}
\pagenumbering{gobble}
\begin{titlepage}
\centering
\vspace*{14mm}
{\Huge\bfseries\color{navy}Síntesis de frecuencia de alta precisión\par}
\vspace{2mm}
{\Huge\bfseries\color{navy}con PLL\_SYS y PIO en RP2040\par}
\vspace{8mm}
{\large\color{graytext}Optimización para 7,074 / 14,074 / 28,074 MHz y corrección de Bresenham\par}
\vspace{8mm}
{\color{blue}\rule{0.82\textwidth}{1pt}}
\vspace{10mm}
{\large Proyecto: \textbf{Firmware PixiePico}\\
Plataforma: RP2040 Zero -- Pico SDK -- PIO\\
Fecha: 1 de septiembre de 2026\par}
Autor: Dr. Pedro E. Colla (LU7DZ)
\vspace{17mm}
\begin{minipage}{0.91\textwidth}
{\LARGE\bfseries\color{navy}Resumen del problema\par}
\vspace{3mm}
\begin{technicalbox}
Para una salida cuadrada simple utilizando una arquitectura rp2040 requiere dos estados lógicos por período.
Con el divisor PIO en formato \textbf{16.8}, la frecuencia ideal de un programa
de dos estados es
\[
 f_{\mathrm{out}}=\frac{128f_{\mathrm{sys}}}{N},
\]
donde \(N\) es entero y representa 256 veces el divisor completo. La búsqueda
exhaustiva sobre todos los PLL\_SYS realmente alcanzables entre 125 y 250 MHz
no logra un error estático menor que 1 Hz en los tres ejemplos. Los mínimos
son 5,5266 Hz, 25,5548 Hz y 50,4202 Hz. Sin embargo, alternando duraciones
adyacentes con un acumulador de error, el período medio teórico puede hacerse
igual al objetivo.
\end{technicalbox}
La precisión teórica (algorítmica) no equivale a precisión real: para garantizar
\(\pm1\) Hz a 28,074 MHz la referencia de 12 MHz debe estar calibrada mejor
que 0,036 ppm, condición que el reloj de una placa rp2040 no cumple y cuya falta de precisión
dominará el error final aunque la precisión media sea cero.
\end{minipage}
\vfill
\end{titlepage}
\pagenumbering{arabic}

\section{Alcance y supuestos}
Se intenta desarrollar una única salida de onda cuadrada. Un período
completo contiene un estado alto y uno bajo; por ello el factor de estados es
\(S=2\). El reloj de sistema puede someterse a "overclock" pero por consideraciones
prácticas se limita el intervalo al rango 125 a 250 MHz.
Por otra parte, no todas las frecuencia de clock son viables pueso que éstas se logran
mediante divisores tales que puedan producir el  PLL\_SYS del RP2040
deseado a partir de XOSC=\SI{12}{\mega\hertz} y con la ruta usada
por \texttt{set\_sys\_clock\_pll()} y \texttt{check\_sys\_clock\_hz()}.

\subsection{Restricciones del espacio discreto}
\rowcolors{2}{light}{white}
\begin{tabularx}{\textwidth}{>{\bfseries}p{0.24\textwidth}Y}
\toprule
\rowcolor{navy}\color{white}Elemento & \color{white}Restricción usada\\
\midrule
XOSC & \SI{12000000}{\hertz}\\
PLL refdiv & Fijo según \texttt{PLL\_SYS\_REFDIV} del SDK; no es posible asignar
refdiv arbitrarios que \texttt{set\_sys\_clock\_pll()} no aplica.\\
FBDIV & Entero; se enumera 16--320 y se filtra por VCO válido.\\
VCO & 750--1600 MHz.\\
POSTDIV1/2 & 1--7, con \(\mathrm{POSTDIV1}\geq\mathrm{POSTDIV2}\).\\
PLL\_SYS & 125--250 MHz y división exacta del VCO.\\
Divisor PIO & \(D=d_{\mathrm{int}}+d_{\mathrm{frac}}/256=N/256\);
\(N\geq256\). El hardware utiliza formato 16.8.\\
Forma de onda & Dos estados PIO por período: alto y bajo.\\
\bottomrule
\end{tabularx}

El SDK busca configuraciones con un VCO válido y postdivisores enteros;
\texttt{check\_sys\_clock\_hz()} sólo acepta una frecuencia si el resultado
del PLL es exacto. Esta validación debe ser tenida en cuenta en la solución
pues caso contrario el estado del sistema corre el riesgo de quedar inestable
y otros subsistemas funcionar incorrectamente, por ejemplo el sistema de soporte
a USB.

\subsection{Hipótesis  y límites}
El máximo alcanzado consistentemente en las placas disponibles fue 250 MHz,
a la cual el hardware funcionó correctamente (para un clock nominal sin "overclock" 
de 125 MHz. 270 MHz funciona en alguna placas y produce fallas de USB en otras.
Es por lo tanto necesario limitar rango de búsqueda a 250 MHz.
Este rango parece ser seguro para las unidades probadas, aunque no está en la 
especificación del rp2040 (que como se dijo trabaja con un clock de 125 MHz).
Los cálculos suponen un XOSC ideal de 12 MHz; la tolerancia real debe medirse y compensarse.

\clearpage
\section{Modelo matemático}
\subsection{Reloj del PLL}
De acuerdo al método estándar propuesto por el  SDK como referencia
\(f_{\mathrm{ref}}=\SI{12}{\mega\hertz}\), un candidato de sistema es:
\begin{technicalbox}
\[
 f_{\mathrm{sys}}=
 \frac{\num{12000000}\,\mathrm{FBDIV}}
 {\mathrm{POSTDIV1}\,\mathrm{POSTDIV2}}.
\]
\end{technicalbox}
Sólo se conserva el candidato si
\(\mathrm{VCO}=\SI{12}{\mega\hertz}\cdot\mathrm{FBDIV}\) está dentro del
rango del PLL, los postdivisores son válidos,
\(f_{\mathrm{sys}}\in[125,250]\) MHz y
\texttt{check\_sys\_clock\_hz()} confirma la misma configuración.

\subsection{Divisor PIO y salida de dos estados}
El divisor 16.8 del PIO se representa mediante el entero \(N=256D\). Para dos
instrucciones o estados equivalentes por período:
\begin{technicalbox}
\[
 D=\frac{N}{256},\qquad
 f_{\mathrm{out}}=\frac{f_{\mathrm{sys}}}{2D}
                  =\frac{128f_{\mathrm{sys}}}{N}.
\]
\end{technicalbox}
Para cada candidato PLL:
\begin{technicalbox}
\[
 N^\star=\frac{128f_{\mathrm{sys}}}{f_{\mathrm{objetivo}}},
 \qquad N_0=\lfloor N^\star\rfloor,\qquad N_1=N_0+1.
\]
\end{technicalbox}
Basta evaluar \(N_0\) y \(N_1\) porque \(f_{\mathrm{out}}\) es monótona en
\(N\). La optimización global escoge la combinación que minimiza
\[
 \left|\frac{128f_{\mathrm{sys}}}{N}-f_{\mathrm{objetivo}}\right|.
\]

\subsection{Pseudocódigo del algoritmo}
\begin{lstlisting}[style=code]
best.error = infinito
para cada FBDIV permitido:
    VCO = 12 MHz * FBDIV
    si VCO no es valido: continuar
    para POSTDIV1=1..7 y POSTDIV2=1..POSTDIV1:
        si VCO no divide exactamente: continuar
        f_sys = VCO/(POSTDIV1*POSTDIV2)
        si f_sys no pertenece a [125,250] MHz: continuar
        si check_sys_clock_hz(f_sys) falla: continuar
        N_ideal = 128*f_sys/f_objetivo
        para N en {floor(N_ideal), ceil(N_ideal)}:
            si N < 256: continuar
            f_out = 128*f_sys/N
            error = f_out - f_objetivo
            conservar si |error| mejora best
\end{lstlisting}

\clearpage
\section{Resultados parciales}
\scriptsize
\rowcolors{2}{light}{white}
\begingroup
\setlength{\tabcolsep}{1.5pt}
\begin{tabularx}{\textwidth}{@{}
  p{0.13\textwidth} p{0.13\textwidth} p{0.21\textwidth}
  p{0.18\textwidth} R{0.19\textwidth} R{0.12\textwidth}@{}}
\toprule
\rowcolor{navy}
\color{white}Objetivo & \color{white}PLL\_SYS &
\color{white}PLL (VCO; FB; P1/P2) &
\color{white}\(N\) y divisor 16.8 &
\color{white}Salida & \color{white}Error\\
\midrule
7,074 MHz & 180,000 MHz & 900 MHz; 75; 5/1 &
\(N=3257\), \(D=12+185/256\) &
\num{7073994.473442} Hz & \num{-5.526558} Hz\\
14,074 MHz & 163,500 MHz & 1308 MHz; 109; 4/2 &
\(N=1487\), \(D=5+207/256\) &
\num{14073974.445192} Hz & \num{-25.554808} Hz\\
28,074 MHz & 130,500 MHz & 1044 MHz; 87; 4/2 &
\(N=595\), \(D=2+83/256\) &
\num{28073949.579832} Hz & \num{-50.420168} Hz\\
\bottomrule
\end{tabularx}
\endgroup
\normalsize

La tabla es el resultado de una búsqueda exhaustiva sobre 246 frecuencias
PLL configurables distintas. Bajo los supuestos definidos no existe otra
combinación estática PLL/divisor 16.8 con menor error.

\subsection{Caso \SI{7074000}{\hertz}}
\[
\mathrm{VCO}=\SI{12}{\mega\hertz}\cdot75=\SI{900}{\mega\hertz},
\qquad f_{\mathrm{sys}}=\frac{900}{5\cdot1}=\SI{180}{\mega\hertz}.
\]
\[
N^\star=\frac{128\cdot\num{180000000}}{\num{7074000}}
=\frac{1280000}{393}=3256.997455470738.
\]
\[
\begin{aligned}
N_0&=3256,&D_0&=12+\frac{184}{256},&
f_0&=\SI{7076167.076167076}{\hertz},&
e_0&=+\SI{2167.076167076}{\hertz},\\
N_1&=3257,&D_1&=12+\frac{185}{256},&
f_1&=\SI{7073994.473441818}{\hertz},&
e_1&=-\SI{5.526558182}{\hertz}.
\end{aligned}
\]

\subsection{Caso \SI{14074000}{\hertz}}
\[
\mathrm{VCO}=\SI{12}{\mega\hertz}\cdot109=\SI{1308}{\mega\hertz},
\qquad f_{\mathrm{sys}}=\frac{1308}{4\cdot2}=\SI{163.5}{\mega\hertz}.
\]
\[
N^\star=\frac{128\cdot\num{163500000}}{\num{14074000}}
=\frac{10464000}{7037}=1486.997299985789.
\]
\[
\begin{aligned}
N_0&=1486,&D_0&=5+\frac{206}{256},&
f_0&=\SI{14083445.491251683}{\hertz},&
e_0&=+\SI{9445.491251682}{\hertz},\\
N_1&=1487,&D_1&=5+\frac{207}{256},&
f_1&=\SI{14073974.445191661}{\hertz},&
e_1&=-\SI{25.554808339}{\hertz}.
\end{aligned}
\]

\subsection{Caso \SI{28074000}{\hertz}}
\[
\mathrm{VCO}=\SI{12}{\mega\hertz}\cdot87=\SI{1044}{\mega\hertz},
\qquad f_{\mathrm{sys}}=\frac{1044}{4\cdot2}=\SI{130.5}{\mega\hertz}.
\]
\[
N^\star=\frac{128\cdot\num{130500000}}{\num{28074000}}
=\frac{2784000}{4679}=594.998931395597.
\]
\[
\begin{aligned}
N_0&=594,&D_0&=2+\frac{82}{256},&
f_0&=\SI{28121212.121212121}{\hertz},&
e_0&=+\SI{47212.121212121}{\hertz},\\
N_1&=595,&D_1&=2+\frac{83}{256},&
f_1&=\SI{28073949.579831932}{\hertz},&
e_1&=-\SI{50.420168067}{\hertz}.
\end{aligned}
\]

\clearpage
\section{Corrección temporal tipo Bresenham}
Es claro que la tolerancia del método es adecuada para algunas aplicaciones
(por ejemplo CW) pero no para otras (por ejemplo FT8).
Se busca un método de compensación calculando una frecuencia media a partir
de la utilización del método de Bresenham.
Si se alternan \(N_0=\lfloor N^\star\rfloor\) y \(N_1=N_0+1\), el
parámetro correcto que debe promediarse es \(N\), porque el período es
lineal en \(N\):
\begin{technicalbox}
\[
T(N)=\frac{N}{128f_{\mathrm{sys}}}.
\]
\end{technicalbox}
Sea \(r=N^\star-N_0\). Se utiliza \(N_1\) durante la fracción \(r\) de los
períodos y \(N_0\) durante \(1-r\). Entonces:
\[
\overline N=(1-r)N_0+rN_1=N^\star,\qquad
\overline T=\frac{1}{f_{\mathrm{objetivo}}}.
\]
El error medio teórico es cero. No se deben promediar las frecuencias
\(f(N_0)\) y \(f(N_1)\), porque la magnitud físicamente acumulada es el
período o la fase.

\scriptsize
\rowcolors{2}{light}{white}
\begin{tabularx}{\textwidth}{
  p{0.13\textwidth} p{0.10\textwidth} p{0.10\textwidth}
  Y p{0.15\textwidth} p{0.15\textwidth}}
\toprule
\rowcolor{navy}
\color{white}Objetivo & \color{white}\(N_0\) corto &
\color{white}\(N_1\) largo & \color{white}Uso exacto por patrón &
\color{white}Longitud & \color{white}Repetición\\
\midrule
7,074 MHz & 3256 & 3257 & \(N_1\): 392; \(N_0\): 1 &
393 períodos & 18.000 Hz\\
14,074 MHz & 1486 & 1487 & \(N_1\): 7018; \(N_0\): 19 &
7037 períodos & 2.000 Hz\\
28,074 MHz & 594 & 595 & \(N_1\): 4674; \(N_0\): 5 &
4679 períodos & 6.000 Hz\\
\bottomrule
\end{tabularx}
\normalsize

Para 7,074 MHz,
\[
N^\star=3256+\frac{392}{393}.
\]
La secuencia contiene 392 períodos con \(N=3257\) y uno con \(N=3256\).
Bresenham distribuye el período corto sin acumular más de una unidad de
error del acumulador.

\subsection{Acumulador entero}
\begin{lstlisting}[style=code,language=C]
// r = r_num/r_den; devuelve N0 o N1 para el proximo periodo
uint32_t next_N(fracn_t *s) {
    s->acc += s->r_num;
    if (s->acc >= s->r_den) {
        s->acc -= s->r_den;
        return s->N1;       // periodo ligeramente mas largo
    }
    return s->N0;           // periodo ligeramente mas corto
}
\end{lstlisting}
La secuencia exacta es periódica y genera espurias a la
frecuencia de repetición indicada y sus armónicos. Si el espectro exige
menor concentración tonal, puede emplearse difusión de error de orden
superior o aleatorización controlada, conservando el número de
\(N_0/N_1\) en ventanas largas.

\clearpage
\section{Implementación en RP2040 PIO}
\subsection{Limitación}
\begin{technicalbox}
Una state machine PIO no puede escribir su propio registro \texttt{CLKDIV}.
Por tanto, alternar directamente \(N_0/N_1\) no puede implementarse con
instrucciones PIO puras. Requiere cooperación de CPU/DMA o un generador
temporal PIO rediseñado que alterne duraciones de estado. Cambiar
\texttt{CLKDIV} desde una interrupción por cada período sería inviable a
7--28 MHz y produciría jitter en magnitudes no deseables.
\end{technicalbox}

\subsection{Arquitectura: PIO + DMA + patrón de tiempo}
La alternativa posible es no tratar de reescribir \texttt{CLKDIV} en cada ciclo.
Se fija un reloj de la state machine y se entrega por DMA un flujo de
decisiones corto/largo generado por el acumulador Bresenham. Cada bit
determina si el siguiente semiperíodo consume \(K\) o \(K+1\) ticks de la
state machine. La PIO usa side-set para conmutar el pin y una rama balanceada
para insertar el tick opcional. El solver debe incluir los ciclos de
\texttt{out}/\texttt{jmp}/\texttt{nop} en \(K\); no puede reutilizar sin
cambios la fórmula ideal de dos instrucciones.

Con autopull de 32 bits, incluso 56,148 millones de decisiones de
semiperíodo por segundo requieren aproximadamente 7,02 MB/s de lectura DMA,
una carga razonable desde SRAM. El patrón debe residir en SRAM y usar un
anillo DMA; no debe leerse desde flash mientras otras funciones escriben XIP.

\begin{lstlisting}[style=code]
; Script PIO un bit por semiperiodo
.program bresenham_square
.side_set 1
.wrap_target
    out x, 1       side 1 [BASE_H] ; bit=1 agrega un tick
    jmp !x high_ok side 1
    nop            side 1
high_ok:
    out x, 1       side 0 [BASE_L]
    jmp !x low_ok  side 0
    nop            side 0
low_ok:
.wrap
; autopull=32; DMA alimenta TX FIFO desde un buffer circular
\end{lstlisting}

Este fragmento es por ahora un esqueleto: \texttt{BASE\_H} y
\texttt{BASE\_L}, junto con los ciclos de las ramas, deben resolverse para
cada banda de frecuencia, verificando que ambos caminos tengan exactamente
\(K\) o \(K+1\) ticks. La salida cambia sólo en instrucciones con side-set,
por lo que la fase queda acotada a un tick de la state machine.

\subsection{Alternativa, dither del registro CLKDIV}
También puede escribirse una tabla de \(N_0/N_1\) en
\texttt{PIO\_SMx\_CLKDIV} mediante DMA sincronizado a límites de período.
Conserva un jitter menor, pero exige un mecanismo adicional de pacing y
sincronización porque la state machine no genera por sí sola un DREQ de fin
de período. A 28 MHz, transferir un word por período consume cerca de
112 MB/s y queda demasiado próximo al límite práctico; por eso no es la
opción recomendada salvo que se use compresión por rachas y sincronización.

\clearpage
\section{Solver y flujo de firmware}
\rowcolors{2}{light}{white}
\begin{tabularx}{\textwidth}{p{0.16\textwidth}Y p{0.29\textwidth}}
\toprule
\rowcolor{navy}
\color{white}Etapa & \color{white}Responsabilidad &
\color{white}Fallo seguro\\
\midrule
1. solve & Enumera PLL + divisor; no toca hardware. &
Devuelve \texttt{false} si no hay solución.\\
2. validate & \texttt{check\_sys\_clock\_hz}; límites VCO/PLL/PIO/error. &
Rechaza el candidato antes de aplicarlo.\\
3. apply & Aplica PLL antes de USB y verifica \texttt{clock\_get\_hz}. &
Vuelve a 250 MHz o reinicia.\\
4. USB & \texttt{TUDstart} sólo con reloj estable. &
No reconfigura PLL con USB montado.\\
5. generate & PIO fijo o PIO+DMA Bresenham. &
Detiene la salida ante underrun.\\
\bottomrule
\end{tabularx}

Para un cambio de frecuencia en ejecución, el solver primero comprueba si el
PLL actual permite una solución aceptable cambiando sólo el patrón o divisor.
Si exige otro PLL, se guarda la frecuencia en watchdog scratch y se reinicia;
reprogramar PLL\_SYS después de enumerar dispositivos como parte del 
subsistema TinyUSB (que opera concurrentemente) probablemente genere 
inestabilidades en el firmware.

\subsection{Criterio de selección}
El primer criterio es error medio \(\leq1\) Hz. Entre candidatos que cumplen,
se minimizarán, en este orden: cambio respecto del PLL actual, tasa de
correcciones Bresenham, longitud y tonalidad del patrón, jitter pico y
consumo. Escoger únicamente el menor error estático puede producir una
secuencia de corrección espectralmente peor.

\subsection{Validaciones de ejecución}
\rowcolors{2}{light}{white}
\begin{tabularx}{\textwidth}{p{0.22\textwidth}Y}
\toprule
\rowcolor{navy}
\color{white}Prueba & \color{white}Condición de aceptación\\
\midrule
PLL aplicado & \texttt{clock\_get\_hz(clk\_sys)} coincide con la solución;
\texttt{clk\_usb} permanece en 48 MHz.\\
Divisor & \(N\geq256\); INT/FRAC coinciden con el registro \texttt{CLKDIV}.\\
DMA & Sin underrun durante al menos 10 minutos; contador de fallos igual a cero.\\
Frecuencia & Medición con contador calibrado; media a menos de 1 Hz del objetivo.\\
Espectro & Espurias del patrón dentro del límite definido para el transmisor.\\
USB & CDC/Audio/MSC permanecen enumerados durante arranque y operación.\\
\bottomrule
\end{tabularx}

\clearpage
\section{Precisión absoluta y calibración}
\rowcolors{2}{light}{white}
\begin{tabularx}{0.78\textwidth}{p{0.22\textwidth}p{0.22\textwidth}Y}
\toprule
\rowcolor{navy}
\color{white}Frecuencia & \color{white}1 Hz relativo &
\color{white}Estabilidad mínima de referencia\\
\midrule
7,074 MHz & \(1{,}414\times10^{-7}\) & 0,141363 ppm\\
14,074 MHz & \(7{,}105\times10^{-8}\) & 0,071053 ppm\\
28,074 MHz & \(3{,}562\times10^{-8}\) & 0,035620 ppm\\
\bottomrule
\end{tabularx}

Un cristal de 10 ppm puede introducir aproximadamente 70,74 Hz, 140,74 Hz y
280,74 Hz, respectivamente. Por lo tanto, la especificación física de
\(\pm1\) Hz exige calibración frente a una referencia mejor, compensación por
temperatura o disciplina externa, por ejemplo un GPSDO o TCXO adecuado.
Bresenham elimina cuantización media; no elimina error de referencia, ruido
de fase ni deriva térmica.
Sin embargo en todos los métodos de comunicación objetivo de éste proyecto
importa mas el error relativo que el error absoluto.

\section*{Conclusiones}
\addcontentsline{toc}{section}{Conclusiones}
\begin{technicalbox}
El análisis anterior requiere implementar primero el solver exhaustivo con las
validaciones recomendadas en el PICO-SDK para a continuación
registrar las mejores soluciones estáticas y confirmar los números de esta
nota. 
Finalmente implementar el generador PIO+DMA con acumulador de fase, medir
espurias y sólo entonces integrarlo al firmware USB. Dejar de reconfigurar
PLL\_SYS con TinyUSB montado; aplicar el PLL seleccionado antes de
\texttt{TUDstart()} o reiniciar ante cambios que exijan otro PLL.
\end{technicalbox}

Si bien se estima como necesario una precisión de frecuencia de alrededor de
1 Hz no puede garantizarse físicamente sólo
con búsqueda PLL y divisor PIO. Sí puede lograrse error medio de cuantización
menor que 1 Hz ---incluso cero teórico--- mediante modulación temporal, pero
la precisión absoluta requiere una referencia calibrada y validación
espectral, sin embargo se comentó en la sección anterior que no es tan 
importante en los casos de uso.

\section*{Referencias}
\begin{enumerate}
\item Raspberry Pi Ltd., \emph{RP2040 Datasheet}, secciones de PLL, clocks y PIO.
\item Raspberry Pi Pico SDK, \texttt{hardware\_clocks}:
\texttt{check\_sys\_clock\_hz()}, \texttt{set\_sys\_clock\_pll()} y
\texttt{clock\_get\_hz()}.
\item Raspberry Pi Pico SDK, \texttt{hardware\_pio}: divisor de clock 16.8 y
\texttt{pio\_sm\_set\_clkdiv\_int\_frac8()}.
\item TinyUSB / Pico SDK: mantener \texttt{clk\_usb} independiente y aplicar
el PLL de sistema antes de enumerar el dispositivo.
\item Código del proyecto Firmware PixiePico, módulos \texttt{quad.c},
\texttt{quad.h} y checkpoint del 31-08-2026.
\end{enumerate}
\end{document}
